% !TEX root = ../main.tex
%=====================================================================
%  参考文献（原文 pp.78--81，编号 [1]--[58]）
%  条目内容按原文照录，作者名与标题保留英文；原文本身残缺者照录并加注。
%=====================================================================

\begin{thebibliography}{99}
\setlength{\itemsep}{2pt}

\bibitem{ref1} M. Atiyah, N. Hitchin, and I. Singer, ``Self-Duality in Four-Dimensional
Riemannian Geometry,'' Proc.\ Royal Soc.\ (London) A \textbf{362} (1978) 425--461.

\bibitem{ref2} M. Atiyah, ``New Invariants Of 3- And 4-Dimensional Manifolds,'' in
\emph{The Mathematical Heritage of Hermann Weyl}, Proc.\ Symp.\ Pure Math.\ \textbf{48},
AMS 1988.

\bibitem{ref3} M. F. Atiyah and L. Jeffrey, ``Topological Lagrangians and cohomology,''
J.\ Geom.\ Phys.\ \textbf{7}, 119 (1990).

\bibitem{ref4} S. Cordes, G. W. Moore and S. Ramgoolam, ``Lectures on 2-d Yang-Mills
theory, equivariant cohomology and topological field theories,'' Nucl.\ Phys.\ Proc.\
Suppl.\ \textbf{41}, 184 (1995) [hep-th/9411210].

\bibitem{ref5} A. Daemi and Y. Xie, ``Sutured Manifolds and Polynomial Invariants from
Higher Rank Bundles,'' arXiv:1701.00571 [math.GT].

\bibitem{ref6} P. Deligne and J. W. Morgan, \emph{Notes on Supersymmetry (following
Joseph Bernstein)}, in \emph{Quantum Fields and Strings: A Course for Mathematicians},
Vol.~1, pp.~41--96, AMS.

\bibitem{ref7} S. K. Donaldson, ``Yang-Mills Invariants of Four-Manifolds,'' in
\emph{Geometry of Low-Dimensional Manifolds: Gauge theory and algebraic surfaces},
S. K. Donaldson and C. B. Thomas eds., London Math.\ Soc.\ Lecture Note Series 150.

\bibitem{ref8} S. Donaldson and P. Kronheimer, \emph{The Geometry of Four-Manifolds},
Oxford 1990.

\bibitem{ref9} S. Donaldson, Bull.\ Amer.\ Math.\ Soc.\ Vol.~\textbf{33} (1996) p.~45.

\bibitem{ref10} J. D. Edelstein, M. Gomez-Reino and M. Marino, ``Blowup formulae in
Donaldson-Witten theory and integrable hierarchies,'' Adv.\ Theor.\ Math.\ Phys.\
\textbf{4}, 503 (2000) [hep-th/0006113].

\bibitem{ref11} Paul M. N. Feehan, Peter B. Kronheimer, Thomas G. Leness, Tomasz S.
Mrowka, ``PU(2) monopoles and a conjecture of Marino, Moore, and Peradze,''
arXiv:math/9812125 [math.DG].

\bibitem{ref12} ``Six lectures on four manifolds.''
% [Tutor Note: 【原稿残缺】原文此条仅有标题，无作者与出处，作者显然尚未补全。照录。]

\bibitem{ref13} M. H. Freedman, A. Kitaev, C. Nayak, J. K. Slingerland, K. Walker,
Z. Wang, ``Universal manifold pairings and positivity,'' arXiv:math/0503054 [math.GT].

\bibitem{ref14} P. Feehan and T. Leness, ``SERIES OF PAPERS.''
% [Tutor Note: 【原稿残缺】原文此条为占位符，作者标注为"系列论文"，未给出具体条目。
%  指 Feehan--Leness 关于 PU(2) 单极子与 Witten 猜想的系列长文（含
%  arXiv:math/9802056, math/9905080 等）。照录原状。]

\bibitem{ref15} D. S. Freed and K. K. Uhlenbeck, \emph{Instantons and Four-Manifolds},
MSRI Publications 1984.

\bibitem{ref16} R. Friedman and J. Morgan, \emph{Smooth Four-Manifolds and Complex
Surfaces}, Springer 1994.

\bibitem{ref17} Gompf and Stipsicz, \emph{4-Manifolds and Kirby Calculus}.
% [Tutor Note: 【原稿残缺】原文此条只有作者姓与书名，无出版社与年份；按上下文应为
%  R. E. Gompf and A. I. Stipsicz, \emph{4-Manifolds and Kirby Calculus},
%  AMS Graduate Studies in Mathematics 20, 1999。]

\bibitem{ref18} L. G\"ottsche, H. Nakajima, and K. Yoshioka, ``Donaldson=Seiberg-Witten
From Mochizuki's Formula and Instanton Counting,'' arXiv:1001.5024.

\bibitem{ref19} Allyn Jackson, ``A Revolution in Mathematics,'' Notices AMS 1995 ??
% [Tutor Note: 【原稿残缺】原文以 ``??'' 结尾，卷期页码未定。照录。]

\bibitem{ref20} D. Kotschick, ``Gauge Theory is Dead! -- Long Live Gauge Theory!''
Notices AMS ???
% [Tutor Note: 【原稿残缺】出处未定，原文以 ``???'' 结尾。照录。]

\bibitem{ref21} P. B. Kronheimer and T. S. Mrowka, ``Embedded surfaces and the structure
of Donaldsons polynomial invariants,'' J.\ Differential Geom.\ \textbf{41} (1995), no.~3,
573--734.

\bibitem{ref22} P. Kronheimer, ``Four-manifold invariants from higher-rank bundles,''
arXiv:math/0407518.

\bibitem{ref23} P. B. Kronheimer, ``Some non-trivial families of symplectic structures,''
On his homepage.

\bibitem{ref24} J. Labastida and M. Marino, ``Topological quantum field theory and four
manifolds,'' (Mathematical physics studies. 25).

\bibitem{ref25} D. Leites, ``Introduction to supermanifolds,'' 1980 Russ.\ Math.\ Surv.\
\textbf{35}, 1.

\bibitem{ref26} A. Losev, N. Nekrasov and S. L. Shatashvili, ``Issues in topological gauge
theory,'' Nucl.\ Phys.\ B \textbf{534}, 549 (1998) [hep-th/9711108].

\bibitem{ref27} T.-J. Li and Ai-Ko Liu, ``Family Seiberg-Witten invariants and wall
crossing formulas,'' arXiv:math/0107211.

\bibitem{ref28} J. M. Maldacena, G. W. Moore and N. Seiberg, ``D-brane instantons and
K theory charges,'' JHEP \textbf{0111}, 062 (2001) doi:10.1088/1126-6708/2001/11/062
[hep-th/0108100].

\bibitem{ref29} M. Marino and G. W. Moore, ``Integrating over the Coulomb branch in N=2
gauge theory,'' Nucl.\ Phys.\ Proc.\ Suppl.\ \textbf{68}, 336 (1998) [hep-th/9712062].

\bibitem{ref30} M. Marino and G. W. Moore, ``Donaldson invariants for nonsimply connected
manifolds,'' Commun.\ Math.\ Phys.\ \textbf{203}, 249 (1999) [hep-th/9804104].

\bibitem{ref31} M. Marino and G. W. Moore, ``The Donaldson-Witten function for gauge groups
of rank larger than one,'' Commun.\ Math.\ Phys.\ \textbf{199}, 25 (1998)
[hep-th/9802185].

\bibitem{ref32} M. Marino and G. W. Moore, ``Three manifold topology and the
Donaldson-Witten partition function,'' Nucl.\ Phys.\ B \textbf{547}, 569 (1999)
[hep-th/9811214].

\bibitem{ref33} M. Marino, G. W. Moore and G. Peradze, ``Superconformal invariance and the
geography of four manifolds,'' Commun.\ Math.\ Phys.\ \textbf{205}, 691 (1999)
[hep-th/9812055].

\bibitem{ref34} M. Marino, G. W. Moore and G. Peradze, ``Four manifold geography and
superconformal symmetry,'' math/9812042.

\bibitem{ref35} J. Milnor, ``On simply connected 4-manifolds,'' 1958 Symposium
internacional de topolog\'ia algebraica / International symposium on algebraic topology,
pp.~122--128, Universidad Nacional Aut\'onoma de M\'exico and UNESCO, Mexico City.

\bibitem{ref36} T. Mochizuki, ``A theory of the invariants obtained from the moduli stacks
of stable objects on a smooth polarized surface,'' arXiv:math/0210211 [math.AG].

\bibitem{ref37} G. W. Moore, ``2-D Yang-Mills theory and topological field theory,''
hep-th/9409044.

\bibitem{ref38} G. W. Moore and E. Witten, ``Integration over the u plane in Donaldson
theory,'' Adv.\ Theor.\ Math.\ Phys.\ \textbf{1}, 298 (1998) [hep-th/9709193].

\bibitem{ref39} G. W. Moore, N. Nekrasov and S. Shatashvili, ``Integrating over Higgs
branches,'' Commun.\ Math.\ Phys.\ \textbf{209}, 97 (2000) doi:10.1007/PL00005525
[hep-th/9712241].

\bibitem{ref40} G. W. Moore, N. Nekrasov and S. Shatashvili, ``D particle bound states and
generalized instantons,'' Commun.\ Math.\ Phys.\ \textbf{209}, 77 (2000)
doi:10.1007/s002200050016 [hep-th/9803265].

\bibitem{ref41} G. W. Moore, A. B. Royston and D. Van den Bleeken, ``L2-Kernels Of
Dirac-Type Operators On Monopole Moduli Spaces,'' arXiv:1512.08923 [hep-th].

\bibitem{ref42} J. W. Morgan, \emph{The Seiberg-Witten equations and applications to the
topology of smooth four-manifolds}, Princeton University Press Mathematical Notes, 1996.

\bibitem{ref43} N. A. Nekrasov, ``Seiberg-Witten prepotential from instanton counting,''
Adv.\ Theor.\ Math.\ Phys.\ \textbf{7}, no.~5, 831 (2003)
doi:10.4310/ATMP.2003.v7.n5.a4 [hep-th/0206161].

\bibitem{ref44} N. Nekrasov and S. Shadchin, ``ABCD of instantons,'' Commun.\ Math.\
Phys.\ \textbf{252}, 359 (2004) doi:10.1007/s00220-004-1189-1 [hep-th/0404225].

\bibitem{ref45} N. Nekrasov and V. Pestun, ``Seiberg-Witten geometry of four dimensional
N=2 quiver gauge theories,'' arXiv:1211.2240 [hep-th].

\bibitem{ref46} N. Seiberg and E. Witten, ``Electric - magnetic duality, monopole
condensation, and confinement in N=2 supersymmetric Yang-Mills theory,'' Nucl.\ Phys.\
B \textbf{426}, 19 (1994) [Erratum-ibid.\ B \textbf{430}, 485 (1994)] [hep-th/9407087].

\bibitem{ref47} N. Seiberg and E. Witten, ``Monopoles, duality and chiral symmetry breaking
in N=2 supersymmetric QCD,'' Nucl.\ Phys.\ B \textbf{431}, 484 (1994) [hep-th/9408099].

\bibitem{ref48} A. Scorpan, \emph{The Wild World of Four-Manifolds}, AMS.

\bibitem{ref49} A. D. Shapere and Y. Tachikawa, ``Central charges of N=2 superconformal
field theories in four dimensions,'' JHEP \textbf{0809}, 109 (2008)
doi:10.1088/1126-6708/2008/09/109 [arXiv:0804.1957 [hep-th]].

\bibitem{ref50} ``Will we every classify 4-folds?''
% [Tutor Note: 【原稿残缺】原文此条仅有标题（且原文即有拼写错误 ``every'' 应为
%  ``ever''），无作者与出处。照录。]

\bibitem{ref51} Y. Tachikawa, ``N=2 supersymmetric dynamics for pedestrians,'' Lect.\
Notes Phys.\ \textbf{890} (2014) doi:10.1007/978-3-319-08822-8 [arXiv:1312.2684
[hep-th]].

\bibitem{ref52} P. Teichner and E. Vogt, ``ALL 4-MANIFOLDS HAVE SPIN$^{c}$ STRUCTURES,''
\texttt{http://people.mpim-bonn.mpg.de/teichner/Papers/spin.pdf}.

\bibitem{ref53} E. Witten, ``Topological Quantum Field Theory,'' Commun.\ Math.\ Phys.\
\textbf{117}, 353 (1988).

\bibitem{ref54} E. Witten, ``Supersymmetric Yang-Mills theory on a four manifold,''
J.\ Math.\ Phys.\ \textbf{35}, 5101 (1994) [hep-th/9403195].

\bibitem{ref55} E. Witten, ``Monopoles and four manifolds,'' Math.\ Res.\ Lett.\
\textbf{1}, 769 (1994) [hep-th/9411102].

\bibitem{ref56} E. Witten, ``On S duality in Abelian gauge theory,'' Selecta Math.\
\textbf{1}, 383 (1995) [hep-th/9505186].

\bibitem{ref57} E. Witten, ``From Superconductors and Four-Manifolds to Weak
Interactions,'' Bull.\ Amer.\ Math.\ Soc.\ Vol.~\textbf{44} (2007) 361.

\bibitem{ref58} E. Witten, ``Notes on Supermanifolds and Integration,'' arXiv:1209.2199.

\end{thebibliography}
